\documentclass{article}

\usepackage{../../Style}

\title{Lie Algebras and their Representations}
\author{Shing Tak Lam}
\date{Michaelmas 2023\footnote{Based on lectures by Stuart Martin. Last updated \today.}}

\makeindex

\begin{document}

\maketitle

\tableofcontents

\lec

\section{Introduction}

A \emph{Lie group} is fundamentally a group, which also a (smooth) manifold. For example, \(\GL_n, \SL_n, \SO_n, \Sp_{2n}\).

\begin{example}
    A prototypical example of a Lie group is the circle group \(S^1\).
\end{example}

Let \(G\) be a Lie group. Then the \emph{Lie algebra}\index{Lie algebra!as tangent space} of \(G\) is the tangent space at the identity \(e\) of \(G\). That is, 
\[\mfg = \T_eG\]
\(\mfg\) is a vector space, with additional structure, which we will see later.

By taking a derivative, we turn the conjugation map
\begin{align*}
    G &\to \Aut(G) \\
    g &\mapsto g(\cdot)g^{-1}
\end{align*}
into a map
\begin{align*}
    \ad : \mfg \to \End(\mfg)
\end{align*}
called the \emph{adjoint}\index{adjoint}. This gives a bilinear map
\begin{align*}
    [\cdot, \cdot] : \mfg \times \mfg &\to \mfg \\
    [x, y] &= \ad(x)(y)
\end{align*}

\begin{example}
    If \(G = \GL_n(\R)\), then we have that \(\mfg = \gl_n(\R) = \Mat_n(\R)\), with
    \[[x, y] = xy - yx\]
\end{example}
What do Lie algebras tell us about the structure of the Lie group \(G\)? 
\begin{itemize}
    \item We will define the \emph{root system} of \(\mfg\), and this then tells us about commutator relations in \(G\) (see Carter's book).
    \item We will define the \emph{Weyl group} of \(\mfg\). For example, the Weyl group of \(\gl_n(\C)\) is isomorphic to \(S_n\), there is an embedding \(S_n \inj \GL_n(\C)\), vie permutation matrices. Let \(B\) denote the Borel subgroup of upper triangular matrices in \(\GL_n(\C)\), then there exists a \emph{Bruhat decomposition}
    \[G = \bigsqcup_{w \in S_n}BwB\]
\end{itemize}
Lie algebras also give us information of the representation theory of \(G\) and \(\mfg\). For example, there exists a bijection
\[\{\text{finite dimensional }\C\text{-representations of }\SL_n(\C)\} \bij \{\text{finite dimensional }\C\text{-representations of }\sl_n(\C)\}\]

Moreover, we can describe the right hand side completely.

In addition, Lie algebras have applications in Algebraic Geometry, for example, we can use Lie algebras to build families of surfaces, of equivalenrly, algebraic curves (See book by Slodovy).

We will define the \emph{Dynkin diagrams} of \emph{semisimple Lie algebras}. For example,
% https://q.uiver.app/#q=WzAsNyxbMCwwLCJcXGJ1bGxldCJdLFsxLDAsIlxcYnVsbGV0Il0sWzIsMCwiXFxidWxsZXQiXSxbMiwxLCJcXGJ1bGxldCJdLFszLDAsIlxcYnVsbGV0Il0sWzQsMCwiXFxidWxsZXQiXSxbNSwwLCJcXGJ1bGxldCJdLFswLDEsIiIsMCx7InN0eWxlIjp7ImhlYWQiOnsibmFtZSI6Im5vbmUifX19XSxbMSwyLCIiLDAseyJzdHlsZSI6eyJoZWFkIjp7Im5hbWUiOiJub25lIn19fV0sWzIsMywiIiwwLHsic3R5bGUiOnsiaGVhZCI6eyJuYW1lIjoibm9uZSJ9fX1dLFsyLDQsIiIsMCx7InN0eWxlIjp7ImhlYWQiOnsibmFtZSI6Im5vbmUifX19XSxbNCw1LCIiLDAseyJzdHlsZSI6eyJoZWFkIjp7Im5hbWUiOiJub25lIn19fV0sWzUsNiwiIiwwLHsic3R5bGUiOnsiaGVhZCI6eyJuYW1lIjoibm9uZSJ9fX1dXQ==
\[\begin{tikzcd}
	\bullet & \bullet & \bullet & \bullet & \bullet & \bullet \\
	&& \bullet
	\arrow[no head, from=1-1, to=1-2]
	\arrow[no head, from=1-2, to=1-3]
	\arrow[no head, from=1-3, to=2-3]
	\arrow[no head, from=1-3, to=1-4]
	\arrow[no head, from=1-4, to=1-5]
	\arrow[no head, from=1-5, to=1-6]
\end{tikzcd}\]
is a Dynkin diagram of type \(E_7\), and understanding the Dynkin diagram tells us about the singularities on the surfaces.

Moreover, Lie algebras also have applications in number theory, root systems/Weyl groups give the structure of groups over \(\Q_p\), the \(p\)-adic integers (see paper on Moodle). Local Langlands correspondence predicts a relationship
\[\{\text{Galois theory of local fields}\}\bij\{\text{complex Lie theory}\}\]

Finally, there any other applications, for example algebraic groups, quantum groups, theoretical physics, quantum mechanics.

\section{Basic definitions and examples}

Let \(k\) be a field. Most of the time, \(k = \C\), but not always. We will sometimes point out how things can go wrong in characteristic \(p\).

\begin{definition}
    [Lie algebra]
    A \emph{Lie algebra}\index{Lie algebra} over \(k\) is a vector space \(\mfg\) over \(k\), together with a bilinear pairing
    \begin{align*}
        [\cdot, \cdot] : \mfg \times \mfg \to \mfg
    \end{align*}
    satisfying
    \begin{enumerate}
        \item \([xx] = 0\) for all \(x \in \mfg\),
        \item the Jacobi identity
        \[[x[yz]] + [y[zx]] + [z[xy]] = 0\]
    \end{enumerate}
\end{definition}

\begin{notation}
    Note that when clear, we will write \([xy] := [x, y]\).
\end{notation}

\begin{remark}
    In particular, we have antisymmetry, i.e.
    \[[xy] = -[yx]\]
\end{remark}

\begin{definition}
    [(Lie) subalgebra]
    A \(k\)-vector subspace \(\mfh\) of \(\mfg\) is a \emph{(Lie) subalgebra}\index{subalgebra@(Lie) subalgebra} if \(\mfh\) is closed under the Lie bracket of \(\mfg\). That is, for all \(x, y \in \mfh\), \([xy] \in \mfh\).
\end{definition}

\begin{example}
    Let \(V\) be a finite dimensional \(k\)-vector space, then
    \begin{enumerate}
        \item Let \(\gl(V) = \End(V)\), with \([xy] = xy-yx\). If we choose a basis for \(V\), then we can identify \(\gl(V)\) with \(\Mat_n(k)\). In this case, we will write \(\gl_n\) or \(\gl_n(k)\)\index{\(\gl_n\)}.
        \item Let
        \[\sl(V) = \left\{x \in \gl(V) \mid \tr(x) = 0\right\}\]
        This defines a subalgebra of \(\gl(V)\), called the \emph{special linear Lie algebra}, with \(\dim(\sl(V)) = \dim(V)^2 - 1\). The standard basis of \(\sl(V)\) is given by
        \begin{align*}
            E_{i,j} \text{ for }i \ne j \\
            E_{i, i} - E_{i+1, i+1}
        \end{align*}
        We will often write \(\sl_n\)\index{\(\sl_n\)} for this Lie algebra.       
    \end{enumerate}
\end{example}

\lec

\begin{example}
    [continued]
    \begin{enumerate}
        \setcounter{enumi}{2}
        \item Suppose \(\Char(k) \ne 2\), and suppose \(V\) is endowed with a symmetric (non-degenerate?) bilinear form
        \[\inner{\cdot, \cdot} : V \times V \to k\]
        Then define
        \[\so(V) = \left\{x \in \gl(V) \mid \inner{xv, w} + \inner{v, xw} = 0 \text{ for all }v, w \in V\right\}\]
        If \(M \in \GL(V)\) is such that \(\inner{v, w} = v^\T M w\), so
        \[\so(V) = \left\{x \mid M x + x^\T M = 0\right\}\]
        We usually take
        \[M = \begin{cases}
            \begin{pmatrix}
                0 & I_\ell \\
                I_\ell & 0
            \end{pmatrix} & n = 2\ell \\
            \begin{pmatrix}
                1 & 0 & 0 \\
                0 & 0 & I_\ell \\
                0 & I_\ell & 0
            \end{pmatrix} & n = 2\ell + 1
        \end{cases}\]
        These are called the \emph{orthogonal Lie algebra}\index{\(\so_n\)}, denoted \(\so_n\).
        \begin{remark}
            \label{rem:sl2c}
            In the case \(n = 2\), let
            \[e = \begin{pmatrix}
                0 & 1 \\ 0 & 0
            \end{pmatrix} \quad f = \begin{pmatrix}
                0 & 0 \\
                1 & 0
            \end{pmatrix} \quad h = \begin{pmatrix}
                1 & 0 \\ 0 & -1
            \end{pmatrix}\]
            viewed as matrices (in \(\sl_2(\C)\)). Indeed, this is the standard basis of \(\sl_2(\C)\). Note
            \begin{align*}
                [ef]&= h \\
                [he] &= 2e \\
                [hf] &= -2f
            \end{align*}
            We'll see that (in some sense) the structure of all semisimple Lie algebras comes from \(\sl_2(\C)\).
        \end{remark}
        \item Again suppose \(\Char(k) \ne 2\). Now suppose \(V\) is endowed with a non-degenerate skew-symmetric (or alternating) bilinear form \(\inner{\cdot, \cdot}\), then
        \[\sp(V) = \left\{x \in \gl(V) \mid \inner{xv, w} + \inner{v, xw} = 0\right\}\]
        In coordinates, we take the form \(\inner{\cdot, \cdot}\) to be the skew-symmetric form given by
        \[M = \begin{pmatrix}
            0 & I_\ell \\
            -I_\ell & 0
        \end{pmatrix}\]
        where \(n = 2\ell\). This is called the \emph{symplectic Lie algebra}\index{\(\sp_n\)}, denoted \(\sp_n\).
    \end{enumerate}
\end{example}

\underline{Exercise:} Check that \(\so_n\) and \(\sp_n\) are Lie subalgebras of \(\gl_n\).

\begin{example}
    [continued]
    \begin{enumerate}
        \setcounter{enumi}{4}
        \item Any vector space \(V\) is a Lie algebra with \([vw] = 0\) for all \(v, w\). We call such Lie algebras \emph{abelian}\index{abelian Lie algebra}. It is named like this, since for \emph{linear Lie algebras}\index{linear Lie algebra}, that is, any subalgebra of \(\gl(V)\) where \(V\) is finite dimensional, \([xy] = xy - yx = 0\) is true if and only if \(x\) and \(y\) commute.
        \item \(\mfb_n\) is the \emph{Borel algebra}\index{\(\mfb_n\)} of upper triangular matrices
        \[\begin{pmatrix}
            * & \cdots & * \\
            & \ddots & \vdots \\
            0& & *
        \end{pmatrix}\]
        This is the Lie algebra associated to the Borel subgroup of upper triangular invertible matrices.
        \item \(\mfn_n\)\index{\(\mfn_n\)} is the Lie algebra of strictly upper triangular matrices,
        \[\begin{pmatrix}
            0 & \cdots & * \\
            & \ddots & \vdots \\
            0& & 0
        \end{pmatrix}\]
        \(\mfn\) stands for \emph{nilpotent}, see section 7.
    \end{enumerate}
\end{example}

\section{Basic representation theory}

\begin{definition}
    [homomorphism, isomorphism]
    A linear map \(\varphi : \mfg \to \mfh\) between two Lie algebras is a \emph{homomorphism}\index{homomorphism} if
    \[\varphi([xy]) = [\varphi(x),\varphi(y)]\]
    for all \(x, y \in \mfg\). If \(\varphi\) is a linear isomorphism, we call \(\varphi\) an \emph{isomorphism of Lie algebras}\index{isomorphism!of Lie algebras}.
\end{definition}

\begin{definition}
    [representation]
    A \emph{representation}\index{representation} of \(\mfg\) is a Lie algebra homomorphism
    \[\varphi : \mfg \to \gl(V)\]
    for some vector space \(V\).
\end{definition}

\begin{notation}
    We also call \(V\) itself a \emph{representation}, or a \(\mfg\)-module. We write \(\mfg \actson V\), and say \(\mfg\) acts on \(V\). We will write
    \[x \vdot v = xv := \varphi(x)(v)\]

    The \emph{dimension}\index{dimension} of the representation is the dimension of \(V\).
\end{notation}

\begin{example}
    \begin{enumerate}
        \item Let \(\dim(V) = 1\), then \(\mfg \actson V\) via \(xv = 0\). This is called the \emph{trivial representation}\index{trivial representation}.
        \item \(\mfg\) is a subalgebra of \(\gl(V)\), then the natural inclusion \(\mfg \inj \gl(V)\) is called the \emph{defining representation}\index{defining representation}.
        \item Let \(x \in \mfg\), define \(\ad_x : \mfg \to \mfg\) by \(\ad_x(y) = [xy]\). The map
        \begin{align*}
            \ad : \mfg &\to \gl(\mfg) \\
            x & \mapsto \ad_x
        \end{align*}
        is called the \emph{adjoint representation}\index{adjoint representation}.
        \begin{remark}
            Recall \(e, h, f\) from \cref{rem:sl2c}, the adjoint representation of \(\sl_2(\C)\) has matrices
            \[\ad_h = \begin{pmatrix}
                2 \\
                & 0 \\
                & & -2
            \end{pmatrix} \quad \ad_e = \begin{pmatrix}
                & -2 \\
                & & 1 \\
                &
            \end{pmatrix} \quad \ad_f = \begin{pmatrix}
                & \\
                -1 \\
                & 2
            \end{pmatrix}\]
            with respect to the basis \(\{e,h,f\}\).
        \end{remark}
        \item If \(V, W\) are representations of \(\mfg\), then so is their \emph{direct sum}\index{direct sum} \(V \oplus V\), via
        \[x(v, w) = (xv, xw)\]
        \item If \(V\) is a representation of \(\mfg\), then so is the \emph{dual}\index{dual representation} \(V^*\),
        \[(xf)(v) = -f(xv)\]
        for all \(x \in \mfg, f \in V^*, v \in V\).
        \item If \(V, W\) are representations of \(\mfg\), then so is the \emph{homomorphisms}\index{hom representation@\(\Hom(V, W)\) representation} \(\Hom(V, W)\), via
        \[(x\vdot f)(v) = x\vdot f(v) - f(x \vdot v)\]
    \end{enumerate}
\end{example}

\begin{definition}
    [equivariant, isomorphism]
    If \(V, W\) are representations of \(\mfg\), then a linear map \(\varphi : V \to W\) is called \emph{\(\mfg\)-equivariant}\index{equivariant} for all \(x \in \mfg, v \in V\),
    \[x \vdot \varphi(v) = \varphi(x \vdot v)\]

    We say \(V, W\) are \emph{isomorphic}\index{isomorphism!of representations} if there exists a \(\mfg\) equivariant isomorphism (of vector spaces) \(V \to W\).
\end{definition}

\begin{definition}
    [subrepresentation, irreducible] A \emph{subrepresentation}\index{subrepresentation} \(W \le V\) is a subspace \(xw \in W\) for all \(x \in \mfg, w \in W\). A non-zero representation \(V\) is \emph{irreducible}\index{irreducible representation} or \emph{simple}\index{simple representation} if the only subrepresentations of \(V\) are \(0\) and \(V\).
\end{definition}

\underline{Exercise:} The trivial representation is irreducible, and so are the defining and the adjoint representation of \(\sl_2(\C)\).

\printindex

\end{document}